(1) 内部粒子满足 \(m\ddot u_j=K(u_{j+1}-2u_j+u_{j-1})\)，\(j=1,\dots,N\)，边界为 \(u_0=u_{N+1}=0\)。动力学矩阵是对称三对角离散Laplacian，因此有实数正本征频率和正交模态。

(2) 设 \(u_j=A\sin(jq)e^{-i\omega t}\)，固定端给 \(q_n=n\pi/(N+1)\)，\(n=1,\dots,N\)。由离散差分 \(2\cos q-2=-4\sin^2(q/2)\)，得 \(\boxed{\omega_n=2\sqrt{K/m}\sin[n\pi/(2(N+1))]}\)。不同模式满足 \(\sum_{j=1}^N\sin(jq_n)\sin(jq_p)=(N+1)\delta_{np}/2\)。

(3) 无限链以物理波数 \(k=q/a\) 描述，第一Brillouin区 \(|k|\le\pi/a\) 的色散是 \(\omega(k)=2\sqrt{K/m}|\sin(ka/2)|\)。对于 \(0<k<\pi/a\)，群速 \(v_g=a\sqrt{K/m}\cos(ka/2)\)，故 \(v_g\) 从 \(a\sqrt{K/m}\) 下降到零；最大频率为 \(2\sqrt{K/m}\)。

(4) 当 \(ka\ll1\) 时 \(\sin(ka/2)\simeq ka/2\)，于是 \(\omega\simeq c|k|\)、\(c=a\sqrt{K/m}\)，恢复连续波动方程。短波接近布里渊区边界时群速趋零而非随频率无限提高。系统机械能为 \(\sum_j m\dot u_j^2/2+(K/2)\sum_{j=0}^N(u_{j+1}-u_j)^2\)，始终非负且守恒。
